\subsection{Factual Consistency Metrics and Their Evaluation}
Early consistency metrics trained classifiers on synthetic corruptions of reference sentences, as FactCC does \citep{kryscinski2020factcc}. Question answering metrics generate questions from the summary and compare the answers obtained from the summary and from the source: QAGS \citep{wang2020qags}, FEQA \citep{durmus2020feqa} and QuestEval \citep{scialom2021questeval} differ mainly in how questions are generated and answers compared, and QAFactEval tunes each of these components \citep{fabbri2022qafacteval}. Entailment metrics score how strongly the source supports each part of the summary, at the level of dependency arcs in DAE \citep{goyal2021dae} or sentence pairs in SummaC \citep{laban2022summac}; AlignScore trains one alignment function on many related tasks \citep{zha2023alignscore}. Meta-evaluation compares such metrics with human labels. FRANK collected typed error annotations and showed that metric rankings depend on the error types a benchmark contains \citep{pagnoni2021frank}; TRUE standardized the comparison as summary-level ROC AUC across datasets \citep{honovich2022true}; and later analyses found that metric quality varies with the summarizer that produced the errors \citep{tang2023understanding}. Large language models are now used as evaluators too, with scores prompted from a strong model \citep{liu2023geval}, although they still miss many factual edits that humans catch \citep{laban2023summedits}. Protocols for the human side have been revisited as well, from atomic fact checking \citep{min2023factscore} to fine-grained content units that raise inter-annotator agreement \citep{liu2023rose}. Almost all of these comparisons use news summaries of sources a few hundred words long, and they report either system-level correlations or summary-level discrimination. We report both, because the two answer different questions: whether a metric can choose between systems, and whether it can flag an individual summary.

\subsection{Long-Document Summarization and Its Evaluation}
Long-document summarization datasets include scientific articles with their abstracts \citep{cohan2018discourse}, government reports \citep{huang2021govreport} and book chapters \citep{kryscinski2022booksum}. Their sources exceed the input length of many metric backbones, so metrics either truncate the source or split it into chunks and aggregate the chunk scores. Both choices can change a verdict: truncation hides the evidence for a statement drawn from late in the document, and aggregation over many chunks can find support for a statement that no single passage makes. SCALE chunks long sources for an entailment model and reports that chunk size strongly affects inconsistency detection \citep{lattimer2023scale}. On the human side, LongEval reports low agreement for faithfulness judgements of long summaries and proposes finer-grained annotation units \citep{krishna2023longeval}. Extractive summaries are not faithful by construction either: copied sentences can mislead once their context is removed \citep{zhang2023extractive}. Our study sits between these lines of work: we take two established news metrics and ask, with an ordinary binary human judgement, how their agreement with annotators changes as documents get longer.
